% Figura (Problema 1.11): tres planos ortogonales cargados; campo por superposición.
\begin{figure}[ht]
\centering
\tdplotsetmaincoords{62}{128}
\begin{minipage}[t]{0.48\textwidth}\centering
\begin{tikzpicture}[tdplot_main_coords, scale=0.88, baseline=(current bounding box.north)]
  \def\L{2.3}
  \fill[azulclaro, opacity=0.30] (0,-\L,-\L) -- (0,\L,-\L) -- (0,\L,\L) -- (0,-\L,\L) -- cycle;   % x=0
  \fill[verdeteo, opacity=0.18] (-\L,0,-\L) -- (\L,0,-\L) -- (\L,0,\L) -- (-\L,0,\L) -- cycle;   % y=0
  \fill[acento, opacity=0.20] (-\L,-\L,0) -- (\L,-\L,0) -- (\L,\L,0) -- (-\L,\L,0) -- cycle;      % z=0
  \draw[azulacero, line width=0.4pt] (0,-\L,-\L) -- (0,\L,-\L) -- (0,\L,\L) -- (0,-\L,\L) -- cycle;
  \draw[verdeteo, line width=0.4pt] (-\L,0,-\L) -- (\L,0,-\L) -- (\L,0,\L) -- (-\L,0,\L) -- cycle;
  \draw[acento, line width=0.4pt] (-\L,-\L,0) -- (\L,-\L,0) -- (\L,\L,0) -- (-\L,\L,0) -- cycle;
  \draw[eje] (0,0,0) -- (3.0,0,0) node[below left] {\small $x$};
  \draw[eje] (0,0,0) -- (0,3.0,0) node[right] {\small $y$};
  \draw[eje] (0,0,0) -- (0,0,3.0) node[above] {\small $z$};
  \node[font=\scriptsize, text=azulprof] at (0,-1.6,1.9) {$\sigma_a$};
  \node[font=\scriptsize, text=verdeteo] at (-1.9,0,1.9) {$\sigma_b$};
  \node[font=\scriptsize, text=acento] at (1.9,1.9,0) {$\sigma_c$};
\end{tikzpicture}\\[2pt]{\small\scshape (a)}
\end{minipage}\hfill
\begin{minipage}[t]{0.48\textwidth}\centering
\begin{tikzpicture}[tdplot_main_coords, scale=0.88, baseline=(current bounding box.north)]
  \draw[eje] (-2.6,0,0) -- (2.8,0,0) node[below left] {\small $x$};
  \draw[eje] (0,-2.6,0) -- (0,2.8,0) node[right] {\small $y$};
  \draw[eje] (0,0,-2.6) -- (0,0,2.8) node[above] {\small $z$};
  % un vector por octante: E = (sgn x, sgn y, sgn z)/2e0 con sigma=1 (módulo constante)
  \foreach \sx in {-1,1}{\foreach \sy in {-1,1}{\foreach \sz in {-1,1}{%
    \pgfmathsetmacro{\ft}{0.55+0.25*(\sx*\sy*\sz)}%
    \mapcolor{\ft}%
    \draw[-{Stealth[length=4.5pt,width=4.5pt]}, line width=1.3pt, mapc] ({1.0*\sx},{1.0*\sy},{1.0*\sz}) -- ++({1.0*\sx},{1.0*\sy},{1.0*\sz});}}}
\end{tikzpicture}\\[2pt]{\small\scshape (b)}
\end{minipage}
\caption{Problema 1.11. (a)~Tres planos ortogonales $x=0$, $y=0$ y $z=0$ con densidades $\sigma_a$, $\sigma_b$ y $\sigma_c$. (b)~Para $\sigma_a=\sigma_b=\sigma_c>0$ el campo, suma
de tres campos uniformes de módulo $\sigma/2\eo$ que cambian de sentido al cruzar su plano, es en cada octante $\frac{\sigma}{2\eo}(\pm\ux\pm\uy\pm\uz)$: constante a trozos y
siempre saliente.}
\label{fig:sem-planos}
\end{figure}
